\section*{Chapter \ref{ch:hf:index-and-basket-arbitrage} --- Index and Basket Arbitrage}

\begin{solution}{exo:hf:index-and-basket-arbitrage:1}
$5\,000\,e^{(0.05-0.015)\times0.25}=5\,043.94$; the band is 5\,043.44 to 5\,044.45.
\end{solution}

\begin{solution}{exo:hf:index-and-basket-arbitrage:2}
Ten micros per E-mini. $1\,000/(5\times150)=1.33$ dollar contracts per Osaka contract.
\end{solution}

\begin{solution}{exo:hf:index-and-basket-arbitrage:3}
$e^{0.002\times0.25}-1\approx5.0$ basis points.
\end{solution}

\begin{solution}{exo:hf:index-and-basket-arbitrage:4}
$e^{-200/400}=0.61$ and $e^{-600/400}=0.22$.
\end{solution}

\begin{solution}{exo:hf:index-and-basket-arbitrage:5}
Adding names past the edge's peak lowers the edge slowly but the tracking noise quickly (from 0.92 to 0.44 basis point between five and fifteen
names), so the ratio keeps rising until the loss of capture and the small names' spreads dominate.
\end{solution}

\begin{solution}{exo:hf:index-and-basket-arbitrage:6}
The first legs are the most likely to arrive in time; putting the largest weights there maximises the weight captured.
\end{solution}

\begin{solution}{exo:hf:index-and-basket-arbitrage:7}
Mean 150 microseconds: the edge peaks at two names and the edge per unit of risk at three. Mean 1\,000 microseconds: ten and forty. The slower the
stocks follow, the more legs arrive in time and the bigger the best basket.
\end{solution}

\begin{solution}{exo:hf:index-and-basket-arbitrage:8}
The legs are not sent at once: in the chapter's simulation the fiftieth leg arrives a millisecond after the future, when most stocks' quotes have
already moved. A backtest that fills every leg at the displayed price overstates the capture of large baskets; the full basket's edge was 0.07 basis
point, not the 1.35 an instantaneous fill would give.
\end{solution}

\begin{solution}{pb:hf:index-and-basket-arbitrage:1}
\begin{enumerate}
\item $F=Se^{(r-q)T}$, and $F(1\pm c)$ with $c$ the costs.
\item The future reacts first and the stocks' quotes follow with delays; arbitrageurs and the stocks' own market makers close the gap.
\item Liquidity and the absolute basis forecast each other, with two-way Granger causality at short horizons, over some 3\,000 days.
\item The slow trade holds to expiry and bears the dividend forecast; the fast one is closed in seconds.
\item See \cref{def:hf:index-and-basket-arbitrage:partial}; $(h-w)^\top\Sigma(h-w)$.
\item Greedily, each name the one that most reduces the tracking variance, with weights solving $\Sigma_{SS}h_S=(\Sigma w)_S$.
\item See \cref{def:hf:index-and-basket-arbitrage:legging}; $e^{-j\delta/\tau}$.
\item It arrives later and it is a smaller name with a wider spread.
\item 1 name: 1.44, 0.93, 0.51; 5: 2.13, 1.21, 0.92; 15: 2.10, 1.39, 0.71; 50: 1.72, 1.65, 0.07 (basis points).
\item Edge at 5 names, edge per unit of risk at 15.
\item The edge peaks at 10 names (1.36) and the ratio at all 50.
\item The best basket is 2 names at 150 microseconds and 10 at a millisecond.
\item At 20 microseconds a leg, net edge 0.51, 0.92, 0.88, 0.71, 0.32 and 0.07 basis point for 1, 5, 10, 15, 30 and 50 names: five names maximise
the edge and fifteen the edge per unit of risk.
\item See \cref{def:hf:index-and-basket-arbitrage:cross}; the ratio of multipliers times currencies.
\item Micro E-minis, one-tenth of the E-mini, since 6 May 2019; the Nikkei on SGX and CME at \textyen500 a point and in Osaka at \textyen1\,000.
\item The \omterm{def:hf:index-and-basket-arbitrage:partial}{partial basket}: its best size is set by the gateway's speed.
\item It keeps the future and the stocks together, so that hedgers in either pay a \omterm{def:hf:fair-value:fair}{fair price}.
\item More, smaller names: the best basket would still be a few dozen, with more tracking error left.
\item From the time between the future's move and each stock's first quote change, by stock and time of day.
\item The race between the legs and the stocks' quotes, priced against spreads and tracking error.
\end{enumerate}
\end{solution}

\begin{solution}{iq:hf:index-and-basket-arbitrage:1}
The fair value's inputs (financing, dividends to expiry), the costs of the basket and the future, whether the stocks have already moved (the gap may
be stale data), and the size available.
\iqlookfor{fair value inputs and executable prices.}
\end{solution}

\begin{solution}{iq:hf:index-and-basket-arbitrage:2}
Minimise $(h-w)^\top\Sigma(h-w)$ over $h$ with $h_j=0$ off $S$: the gradient on $S$ is $2(\Sigma(h-w))_S=0$, so $\Sigma_{SS}h_S=(\Sigma w)_S$; the
objective is convex, so this is the minimum.
\iqlookfor{first-order conditions on the subset.}
\end{solution}

\begin{solution}{iq:hf:index-and-basket-arbitrage:3}
Serialisation on one session and one network card, the exchange's per-session throttle, and the gateway's per-order work; pre-build the orders,
spread them across sessions and cards, and send in parallel.
\iqlookfor{pre-computation and parallel sessions.}
\end{solution}

\begin{solution}{iq:hf:index-and-basket-arbitrage:4}
Long the filled stocks against a future sized for the whole basket: a factor exposure and the unfilled names' tracking error. Resize the future
hedge to the filled legs at once, then decide whether to complete or unwind.
\iqlookfor{hedge what is filled, immediately.}
\end{solution}

\begin{solution}{iq:hf:index-and-basket-arbitrage:5}
Each listing tends to lead in its home market's hours: Osaka in the Tokyo day, Singapore and Chicago when Osaka is closed; measure it rather than
assume it (chapter 8).
\iqlookfor{time-varying leadership, measured.}
\end{solution}

\begin{solution}{iq:hf:index-and-basket-arbitrage:6}
Expected capture $\propto\sum_{j=1}^{k}e^{-j\delta/\tau}$ less $ck$; the marginal leg is worth adding while $e^{-k\delta/\tau}>c$ (in capture
units), so $k^\ast=\lfloor(\tau/\delta)\ln(1/c)\rfloor$.
\iqlookfor{the marginal leg.}
\end{solution}
